\documentclass[12pt]{article}
\usepackage[utf8]{inputenc}
\usepackage[T1]{fontenc}
\usepackage[margin=1in]{geometry}
\usepackage{setspace}
\usepackage{xcolor}
\usepackage[colorlinks=true, urlcolor=blue!60!black, linkcolor=blue!60!black]{hyperref}
\setlength{\parindent}{0em}
\setlength{\parskip}{0.6em}
\renewcommand{\baselinestretch}{1.0}
\begin{document}
\begin{center}
{\LARGE \textit{Timbs v.\ Indiana}, 586 U.S.\ 146 (2019)\par}
\vspace{0.5em}
{\large Case Brief --- Two-Page Edition\par}
\end{center}

\section{Memory Jogger}
The Excessive Fines Clause binds the States, including in civil forfeitures of property like Timbs's \$42,000 Land Rover (\href{https://www.courtlistener.com/opinion/4591916/timbs-v-indiana/}{CourtListener}).

\section{Detailed Case Facts}
Tyson Timbs drove his \$42,000 Land Rover, bought with life insurance proceeds, to sell small amounts of heroin to undercover officers. After he pleaded guilty, Indiana sought civil \textit{in rem} forfeiture of the vehicle, worth more than four times his offense's \$10,000 maximum fine.

\section{Procedural History}
The trial court found the forfeiture ``grossly disproportionate''; the Indiana Supreme Court reversed, holding the Clause unincorporated.

\section{Judicial Votes}
Vacated and remanded, 9--0. Ginsburg wrote for eight Justices; Gorsuch also concurred; Thomas concurred in the judgment.

\section{Holding}
The Excessive Fines Clause is incorporated through the Due Process Clause and applies to the States exactly as to the federal government, including to partly punitive civil forfeitures. The Court did not decide whether this forfeiture was excessive.

\section{Analysis of Opinions}
\textbf{Majority:} From Magna Carta through the English Bill of Rights and the state constitutions of 1787, of 1868, and of today, the protection is deeply rooted; Black Code fines and the revenue incentive of fines show why. Incorporation turns on the right, not each application such as \textit{in rem} forfeiture. \textbf{Gorsuch:} the Privileges or Immunities Clause may be the better vehicle. \textbf{Thomas:} substantive due process is a fiction; the right is a privilege of citizenship.

\section{Examples: Future Applications}
\textbf{Same-Side:} forfeiting a \$150,000 home for a minor offense. \textbf{Opposite-Side:} forfeiting proven drug proceeds or contraband. \textbf{Fence-Sitter:} a modest fine that ruins an indigent defendant's livelihood, since ability to pay was left open.

\section{Critique}
Originalists welcome the result but want the Privileges or Immunities Clause; Beth Colgan and others fear \textit{Timbs} is symbolic without an ability-to-pay test.

\section{Key Quotations}
\begin{quote}
``Incorporated Bill of Rights guarantees are `enforced against the States under the Fourteenth Amendment according to the same standards that protect those personal rights against federal encroachment.' Thus, if a Bill of Rights protection is incorporated, there is no daylight between the federal and state conduct it prohibits or requires.''
\end{quote}
\noindent --- Justice Ginsburg, for the Court, 586 U.S.\ at 150.

\end{document}
